% !TEX TS-program = XeLaTeX
% use the following command:
% all document files must be coded in UTF-8
\documentclass[portuguese]{textolivre}
% build HTML with: make4ht -e build.lua -c textolivre.cfg -x -u article "fn-in,svg,pic-align"

\journalname{Texto Livre}
\thevolume{19}
%\thenumber{1} % old template
\theyear{2026}
\receiveddate{\DTMdisplaydate{2026}{7}{1}{-1}} % YYYY MM DD
\accepteddate{\DTMdisplaydate{2026}{7}{20}{-1}}
\publisheddate{\DTMdisplaydate{2026}{7}{24}{-1}}
\corrauthor{Francis Arthuso Paiva}
\articledoi{10.1590/1983-3652.2026.66990}
%\articleid{NNNN} % if the article ID is not the last 5 numbers of its DOI, provide it using \articleid{} commmand 
% list of available sesscions in the journal: articles, dossier, reports, essays, reviews, interviews, editorial
\articlesessionname{dossier}
\runningauthor{Paiva et al.} 
%\editorname{Leonardo Araújo} % old template
\sectioneditorname{Daniervelin Pereira~\orcid{0000-0003-1861-3609}}
\layouteditorname{Leonardo Araujo~\orcid{0000-0003-3884-2177}}

\title{Editorial: Inteligências artificiais e suas interfaces com a vida social, a educação linguística, a multimodalidade e com o discurso}
\othertitle{Editorial: Artificial intelligences and their interfaces with social life, language education, multimodality, and discourse}
% if there is a third language title, add here:
%\othertitle{Artikelvorlage zur Einreichung beim Texto Livre Journal}

\author[1]{Francis Arthuso Paiva~\orcid{0000-0002-9083-3342}\thanks{Email: \href{mailto:paivafrancis@yahoo.com.br}{paivafrancis@yahoo.com.br}}}
\author[2]{Záira Bomfante dos Santos~\orcid{0000-0002-6162-8489}\thanks{Email: \href{mailto:zbomfante@gmail.com}{zbomfante@gmail.com}}}
\author[3]{Fei Victor Lim~\orcid{0000-0003-3046-1011}\thanks{Email: \href{mailto:victor.lim@nie.edu.sg}{victor.lim@nie.edu.sg}}}
\author[4]{Camila Cárdenas Neira~\orcid{0000-0002-1842-7200}\thanks{Email: \href{mailto:camila.cardenas.neira@gmail.com}{camila.cardenas.neira@gmail.com}}}
\affil[1]{Universidade Federal de Minas Gerais, Colégio Técnico, Belo Horizonte, MG, Brasil.}
\affil[2]{Universidade Federal do Espírito Santo, Vitória, ES, Brasil.}
\affil[3]{Nanyang Technological University, Nanyang Ave, Singapore.}
\affil[4]{Universidad Austral de Chile, Valdivia, Los Ríos, Chile.}

\addbibresource{article.bib}
% use biber instead of bibtex
% $ biber article

% used to create dummy text for the template file
\definecolor{dark-gray}{gray}{0.35} % color used to display dummy texts
\usepackage{lipsum}
\SetLipsumParListSurrounders{\colorlet{oldcolor}{.}\color{dark-gray}}{\color{oldcolor}}

% used here only to provide the XeLaTeX and BibTeX logos
\usepackage{hologo}

% if you use multirows in a table, include the multirow package
\usepackage{multirow}

% provides sidewaysfigure environment
\usepackage{rotating}

% CUSTOM EPIGRAPH - BEGIN 
%%% https://tex.stackexchange.com/questions/193178/specific-epigraph-style
\usepackage{epigraph}
\renewcommand\textflush{flushright}
\makeatletter
\newlength\epitextskip
\pretocmd{\@epitext}{\em}{}{}
\apptocmd{\@epitext}{\em}{}{}
\patchcmd{\epigraph}{\@epitext{#1}\\}{\@epitext{#1}\\[\epitextskip]}{}{}
\makeatother
\setlength\epigraphrule{0pt}
\setlength\epitextskip{0.5ex}
\setlength\epigraphwidth{.7\textwidth}
% CUSTOM EPIGRAPH - END

% to use IPA symbols in unicode add
%\usepackage{fontspec}
%\newfontfamily\ipafont{CMU Serif}
%\newcommand{\ipa}[1]{{\ipafont #1}}
% and in the text you may use the \ipa{...} command passing the symbols in unicode

% LANGUAGE - BEGIN
% ARABIC
% for languages that use special fonts, you must provide the typeface that will be used
% \setotherlanguage{arabic}
% \newfontfamily\arabicfont[Script=Arabic]{Amiri}
% \newfontfamily\arabicfontsf[Script=Arabic]{Amiri}
% \newfontfamily\arabicfonttt[Script=Arabic]{Amiri}
%
% in the article, to add arabic text use: \textlang{arabic}{ ... }
%
% RUSSIAN
% for russian text we also need to define fonts with support for Cyrillic script
% \usepackage{fontspec}
% \setotherlanguage{russian}
% \newfontfamily\cyrillicfont{Times New Roman}
% \newfontfamily\cyrillicfontsf{Times New Roman}[Script=Cyrillic]
% \newfontfamily\cyrillicfonttt{Times New Roman}[Script=Cyrillic]
%
% in the text use \begin{russian} ... \end{russian}
% LANGUAGE - END

% EMOJIS - BEGIN
% to use emoticons in your manuscript
% https://stackoverflow.com/questions/190145/how-to-insert-emoticons-in-latex/57076064
% using font Symbola, which has full support
% the font may be downloaded at:
% https://dn-works.com/ufas/
% add to preamble:
% \newfontfamily\Symbola{Symbola}
% in the text use:
% {\Symbola }
% EMOJIS - END

% LABEL REFERENCE TO DESCRIPTIVE LIST - BEGIN
% reference itens in a descriptive list using their labels instead of numbers
% insert the code below in the preambule:
%\makeatletter
%\let\orgdescriptionlabel\descriptionlabel
%\renewcommand*{\descriptionlabel}[1]{%
%  \let\orglabel\label
%  \let\label\@gobble
%  \phantomsection
%  \edef\@currentlabel{#1\unskip}%
%  \let\label\orglabel
%  \orgdescriptionlabel{#1}%
%}
%\makeatother
%
% in your document, use as illustraded here:
%\begin{description}
%  \item[first\label{itm1}] this is only an example;
%  % ...  add more items
%\end{description}
% LABEL REFERENCE TO DESCRIPTIVE LIST - END


% add line numbers for submission
%\usepackage{lineno}
%\linenumbers

\begin{document}
\maketitle

\section{O dossiê}

Inteligências Artificiais Generativa – IAG – e os algoritmos têm estado cada vez mais presentes na nossa vida. Essa onipresença tem gerado tensões, dúvidas quanto ao seu impacto na vida humana, nas relações sociais, na disseminação de informações e desinformações – \textit{fake news} –, que interferem nas nossas ações, representações e, consequentemente, na vida democrática. Por outro lado, movimentos individuais e coletivos constroem resistências ao uso indiscriminado dos algoritmos em substituição ao trabalhador, bem como mantêm vigilância sobre as relações entre os estados nacionais e iniciativa privada na pessoa das big techs.
 
Desde a realização de tarefas simples a complexas como produzir textos, vídeos, falar, conversar, revisar textos escritos entre outras atividades linguageiras, esses sistemas computacionais têm ganhado cada vez mais adesão e popularidade, consequentemente vêm sendo utilizados por pessoas dos mais variados seguimentos sociais nem sempre com consciência crítica ou discursiva. Em razão disso, muitos debates têm sido suscitados na educação devido ao impacto no ensino, na aprendizagem, no trabalho docente, nas questões que envolvem autoria, criatividade, desenvolvimento da consciência crítica, disseminação de preconceitos, domínio das tecnologias, compreensão das linguagens utilizadas e das transposições multimodais possíveis.
 
Considerando a relevância dessa discussão na atualidade e o seu impacto na vida humana, sinalizamos nossa alegria de reunir nesta edição trabalhos que proporcionam uma visão mais abrangente dos desafios e das potencialidades oferecidas pelo uso cada vez mais difundido de tecnologias algorítmicas no campo da linguagem. Buscamos articular uma equipe de organização com pessoas de diferentes países, Fei Victor Lim, de Singapura, Camila Cárdenas, do Chile, Záira e Francis, do Brasil, no esforço colaborativo voltado à mobilização de diferentes pesquisadores, a fim de compor um caleidoscópio de vozes para tratar de uma temática complexa e sensível para os estudos da Linguagem e da Comunicação.
 
Neste número, há 15 artigos e uma seção especial que traz uma resenha, um ensaio e uma entrevista.  Dois trabalhos buscam mapear pesquisas que têm sido desenvolvidas em contextos distintos com destaque para o Brasil, Singapura, Espanha e países da América Latina. Os demais treze artigos abordam questões relacionadas à educação linguística e semiótica, discutindo possibilidades de articulação entre os processos de ensino e aprendizagem de linguagens e as tecnologias de IAG. São pesquisas realizadas em diversos centros do mundo, tais como Brasil, Singapura, Indonésia, Macau, Inglaterra, Espanha, Uruguai, Chile e EUA.

Ainda que sejam resultados incipientes, neles se configuram o propósito de hibridização de esforços entre humanos e máquinas, sem escantear preocupações éticas, tampouco desconsiderar o arcabouço de conhecimento acumulado nas experiências acadêmicas pré-IAG.

\section{Artigos}
  
O artigo \citetitle{SantosPaivaLim2026}, apresenta um panorama das pesquisas acerca de Inteligência Artificial no Brasil e em Singapura. No contexto brasileiro, os resultados demonstram o esforço da grande área da Linguística em promover primeiras discussões a respeito desse novo fenômeno de linguagem. Por sua vez, as pesquisas em Singapura avançam na área da educação linguística e semiótica movidas pelo grupo de pesquisa AI@NIE.

O artigo \citetitle{PozoSanchez2026InteligenciaArtificialEducacion} revela como os estudos espanhóis concentram 71\% de publicações sobre IA e Educação no contexto hispanoablante, demandando maior inserção dos países latino-americanos nesses estudos.

O artigo \citetitle{Schlemmer2026} chega a conclusões sobre a parceria possível entre IAG e humanos na construção discursiva, embora criando tensões nos conceitos de autoria, criatividade, diaologismo, logo a posição do docente como transmissor de conhecimento é mais uma vez questionada. 

O artigo \citetitle{Cavalcante2026IAUniversidade} procura discutir questões relativas a esse novo ethos docente em tempos de IA. Para isso seus autores discutem o uso da IA para a produção de trabalhos acadêmicos por universitários, levantando as questões principais: Qual o impacto do uso de IA que geram textos na escrita acadêmica de estudantes da graduação? Como identificar textos produzidos por IA nos trabalhos acadêmicos? Que estratégias podem ser usadas para criar atividades em disciplinas de graduação que exijam participação ativa dos estudantes, evitando a dependência de ferramentas de IA? São apontados caminhos para a proposição de atividades que sejam desafiadoras e que dificultem a substituição da produção escrita pela solicitação à IA.

Semelhante resultado é apontado pelo artigo \citetitle{Jatoba2026EducacaoLinguisticaIA}, cujo objetivo era examinar, por meio de um estudo de caso qualitativo, como a Inteligência Artificial é integrada ao ensino de Português como Língua Estrangeira em Macau e quais implicações formativas e avaliativas essa integração produz. Os resultados apontam que a formação docente, em diálogo com o multilinguismo e com o letramento crítico em IA, constitui um eixo decisivo para o uso pedagogicamente responsável dessas tecnologias no ensino de línguas estrangeiras.

Do ponto de vista dos estudantes, o artigo \citetitle{BeroizaValenzuelaPerezPena2026SchoolLeadership} apresenta resultados positivos de uso de IAG por estudantes chilenos, demonstrando maior engajamento emocional e produção de vocabulário mais variado durante as atividades de escrita, mas com uma ressalva: os achados indicam que a IAG tende a ser sobreposta às práticas tradicionais, em vez de transformá-las, exigindo uma liderança sensível ao contexto para alinhar as políticas de inovação às realidades das salas de aula.

Em relação ao uso crítico e ético das IAG para evitar seu uso sobreposto às práticas escolares, o artigo \citetitle{ContrerasDeLaLlaveBallesterPardo2026UsoEticoIA} apresenta estudo com resultados positivos para uso de protocolo ético e transparente, que orientou estudantes de graduação em educação primária a revelarem cônscios e criticamente a quantidade e a qualidade do uso que fizeram de IAG em seus trabalhos acadêmicos de escrita. 

O artigo \citetitle{KlossCordovezFernandezBustamante2026EvaluacionLLM} também contribui para a percepção segundo a qual é preciso criar protocolos de observação e de uso de IAG para a produção de trabalho acadêmicos. O estudo comparou a correção de ensaios acadêmicos por humano especialista e por IAG, concluindo que “la evaluación efectuada por los LLM es similar a la del humano experto, especialmente en la dimensión de ajuste al género en la subdimensión de propósito comunicativo. Sin embargo, hay una baja precisión en las relaciones de cohesión y coherencia, así como con el ajuste a las normas de la lengua". O artigo advoga a necessidade de letrar graduandos calouros para o uso intencional, reflexivo e ético das IAG.

Outro estudo com professores de inglês como língua estrangeira da Indonésia apresenta resultado semelhante a respeito do engajamento em atividades de linguagem criado pela IAG. Ele é reportado no artigo \citetitle{ArifKurniawanWicaksana2026DeterminantsGenAI}, que concluiu “the intention to use generative AI emerged as a robust predictor of actual usage behavior, supporting existing literature on technology adoption and emphasizing the importance of learners’ intentions in engaging with new technologies”.

O artigo \citetitle{SilvaVillarroelOkuyama2026} apresenta um produto, denominado INATICON, que visa a criar uso reflexivo das IAG, ao promover a escalada que vai do seu uso acrítico ao uso responsável pelos estudantes.

O artigo \citetitle{MorenoFernandez2026LenguajeEstereotiposEdadistas} intencionou analisar os preconceitos de idade gerados por diferentes IAG. De acordo com seus resultados, “em termos discursivos, todos os modelos reproduziram estereótipos discriminatórios relacionados à idade, embora o Gemini e o DeepSeek oferecessem representações mais equilibradas e positivas da velhice”.

O artigo \citetitle{GambettaAbellaAbdalaHuertasRehermann2026} apresenta estudo comparativo entre a agência humana e a agência de IAG para realização de análises do sistema de valoração da análise do discurso crítica de corpus formados por tuits. Os resultados da análise comparativa demonstraram que “la herramienta mostró utilidad en las tareas técnicas; sin embargo, evidenció limitaciones interpretativas tales como errores en la clasificación de los posicionamientos de los usuarios, identificación parcial de los ejes temáticos y aplicación restringida del Sistema de la Valoración”.

O artigo \citetitle{DiamantopoulouOrevik2026MultimodalAgency} oferece resultados relevantes para a análise semiótica social multimodal ao propor novos olhares para os conceitos de agência e design relacionados às produções de IAG. Ao analisar designs imagéticos do Duolingo para aprendizagem de línguas, as autoras do London College propõem que a agência levada a cabo por IA é tanto quanto delegada por agentes humanos. 

Observado essa tensão entre agência humana e mediada por algoritmos, o artigo \citetitle{Barros2026VideosYouTube}, com base na análise do discurso crítica e multimodalidade, analisa vídeo do Youtube destinado ao público infantil, concluindo que as interações mediadas por algoritmos tendem a “condicionar o conteúdo produzido na plataforma, porque a affordance do Youtube age como coautora dos discursos nela veiculados”.

Outras tensões são reportadas pelo artigo \citetitle{YouEtAl2026PromptingAIFeedback}, cujos resultados vêm de teste de uso de prompts diversos por estudantes de português como L2 em Macau realizando revisões textuais auxiliados por IAG. Baseados na engenharia de prompt, os resultados apontam para a necessidade de intervenções pedagógicas, para melhorar a qualidade dos prompts criados por estudantes.

\section{Seção especial}

A seção especial do dossiê traz a \citetitle{Barbosa2026ResenhaTerritoriosLetramento}, obra publicada em 2024. Nas palavras da resenhista Vânia Barbosa, “É preciso formar indivíduos sociais autônomos e aptos a manusearem a pluralidade dos textos que circulam nessa sociedade, que, como o livro organizado por Guilherme Brambila vai mostrar, constituem-se em ‘Territórios do Letramento’ nos quais aqueles indivíduos são agentes atuantes na e por meio da linguagem. O convite a um passeio por esses territórios é o que o leitor do livro irá encontrar desde a sua capa, passando por todos os elementos que o constitui”.

Ana Elisa Ribeiro nos brinda com sua escrita sempre instigadora sobre nossa condição ante as novidades linguageiras. Em \citetitle{Ribeiro2026DefinicoesAtualizadasIA}, Ribeiro discute e reflete sobre as implicações do surgimento das IAG nas escolas e como essa tecnologia pode desconstruir o conceito, muitas vezes monolítico, de erro escolarizado.

Por fim, apresentamos \citetitle{KalantzisEtAl2026InterviewsTalks}, realizada por nós organizadores, que versa sobre Inteligência Artificial e ensino na perspectiva da aprendizagem por design, pedagogia dos multiletramentos e da multimodalidade. As questões foram produzidas com base nos artigos dos referidos professores acerca da proposta deles de transpositional grammar e do seu ambiente de IA chamado Cyber Scholar. A entrevista consiste em 5 perguntas temáticas com suas respectivas respostas e referências. 

Desejamos boas leituras para uma comunicação acadêmica e científica proveitosa!


\printbibliography\label{sec-bib}
% if the text is not in Portuguese, it might be necessary to use the code below instead to print the correct ABNT abbreviations [s.n.], [s.l.]
%\begin{portuguese}
%\printbibliography[title={Bibliography}]
%\end{portuguese}


%full list: conceptualization,datacuration,formalanalysis,funding,investigation,methodology,projadm,resources,software,supervision,validation,visualization,writing,review
\begin{contributors}[sec-contributors]
\authorcontribution{Francis Arthuso Paiva}[conceptualization,projadm,supervision,writing,review]
\authorcontribution{Záira Bomfante dos Santos}[conceptualization,projadm,supervision,writing,review]
\authorcontribution{Fei
Victor Lim}[conceptualization,projadm,supervision,writing,review]
\authorcontribution{Camila Cárdenas Neira}[conceptualization,projadm,supervision,writing,review]
\end{contributors}

\begin{dataavailability}
\txtdataavailability{nodata} % options: dataavailable, dataonly, databody, datanotav, nodata
\end{dataavailability}


\end{document}

